\documentclass[11pt]{article}
\usepackage{fontspec}
\usepackage[french,english]{babel}
\usepackage{amsmath}
\usepackage{amssymb}
\usepackage{amsthm}
\usepackage{mathrsfs}
\usepackage{tikz-cd}
\usepackage{enumitem}
\usetikzlibrary{babel}
\newcommand{\printedpage}[1]{\marginpar{\footnotesize p.~#1}}
\newcommand{\msp}[1]{\marginpar{\footnotesize [#1]}}
\newcommand{\Homint}{\operatorname{Hom\,int}}
\setlength{\parindent}{0pt}
\setlength{\parskip}{0.5\baselineskip}
\begin{document}

\printedpage{4}
% head: A. Grothendieck : Dérivateurs, Ch. II, Édition des Universités de Montpellier II et Paris VII
% foot: Version 8 juillet 2001
0-connected (\footnote[4]{This also means that $\alpha$ is $W_0$-coaspheric. For a treatment ab ovo, cf. prop. 3.2, page 43.}). The fact that the $x\backslash C'$ nonempty implies that $\varprojlim\limits_{C} \longrightarrow \varprojlim\limits_{C'}$ is \textit{injective.} For surjectivity, if $\xi' \in \varprojlim\limits_{C'}$ (set-theoretic case), we define $\xi \in \varprojlim\limits_{C}$ by the $\xi(x) \in F(x)$ with the help of % [sic] sujectivité -> surjectivité, lettre r omise
\msp{page 6}
\[ x \xrightarrow{u} \alpha(x') \]
by taking $u^*(\xi(x'))$ (where $\xi(x') \in (\alpha^*(F))(x') = F(\alpha(x')))$, the essential point is to show that the $u^*(\xi(x'))$ depend only on $x$, not on the choice of $(x',u) \in x\backslash C'$. Now for $(x',u)$ variable in $x\backslash C'$, the map
\[
\begin{aligned}
\varphi : \operatorname{Ob}(x\backslash C') &\quad\longrightarrow\quad F(x)\\
(x',u) &\quad\longmapsto\quad u^*(\xi(x'))
\end{aligned}
\]
is such that if there exists an arrow
\[ (x',u) \xrightarrow{u'} (\bar{x}',\bar{u}) \qquad \text{where} \qquad u' : x' \longrightarrow \bar{x}' \]
(thus the diagram
\[
\begin{tikzcd}
 & \alpha(x') \arrow[dd, "\alpha(u')"] \\
x \arrow[ur, "u"] \arrow[dr, "\bar{u}"'] & \\
 & \alpha(\bar{x}')
\end{tikzcd}
\]
is commutative), then $\varphi(x',u) = \varphi(\bar{x}',\bar{u})$. Thus $\varphi$ is constant on the connected components of $x\backslash C'$, thus constant if $x\backslash C'$ is 0-connected, whence $\xi(x)$. We then prove at once that for variable $x$,
\[ \xi \overset{\text{def}}{=} (\xi(x))_{x\in\operatorname{Ob} C} \simeq \varprojlim F, \qquad \text{\texttt{[rather}}\ \ \in \varprojlim F, \text{\texttt{]}} \]
and that $\alpha^*(\xi) = \xi'$, q.e.d.

We can say that for $x_0 \in \operatorname{Ob} C_0$ fixed, if we look at the \textit{constant} presheaf on $x_0\backslash C'$ with value $F(x_0)$, say $\Phi_{x_0}$, the map
\[ (x',u) \longmapsto \xi^*(u') : x_0\backslash C' \longrightarrow F(x_0) \]
can be regarded as a section of this presheaf on $x\backslash C'$, defined by $\xi' \in \varprojlim\limits_{C'} \alpha^*F$, so we have
\[ \varprojlim_{C'} \alpha^*F \xrightarrow{c_{x_0}} \Gamma(x_0\backslash C', \Phi_{x_0}) \qquad (\xleftarrow{\ \sim\ } F(x_0) \text{ if } x_0\backslash C \text{ 0-connected}). \]
If we have \fbox{$x_0 \xrightarrow{v} x_1$} in $C'$, we deduce from this
\[ x_1\backslash C' \xrightarrow{V} x_0\backslash C' \]
(the $x\backslash C'$ are the fibers of a \textit{fibered category} $C\backslash C'$ over $C$ \dots) and we have $F(x_1) \longrightarrow F(x_0)$, whence \fbox{$\Phi_{x_1} \longrightarrow V^*\Phi_{x_0}$} for the corresponding constant sheaves on $x_0\backslash C'$ resp. $x_1\backslash C'$
\printedpage{5}
% head: A. Grothendieck : Dérivateurs, Ch. II, Édition des Universités de Montpellier II et Paris VII
% foot: Version 8 juillet 2001
\msp{page 7}
and we have commutativity in
\[
\begin{tikzcd}[column sep=large]
\varprojlim\limits_{C'} \alpha^*F \arrow[r, "c_{x_1}"] \arrow[d, "c_{x_0}"'] & \Gamma(x_0\backslash C', F(x_0)) \simeq \operatorname{Hom}(\pi_0(x_0\backslash C'), F(x_0)) \arrow[d, "V^*"] & \\
\operatorname{Hom}(\pi_0(x_1\backslash C'), F(x_1)) \simeq \Gamma(x_1\backslash C', F(x_1)) \arrow[r, "\substack{v^*\\ (F(x_1) \longrightarrow F(x_0))}"'] & \Gamma(x_1\backslash C', F(x_0)) \simeq \operatorname{Hom}(\pi_0(x_1\backslash C'), F(x_0)) & \\
\pi_0(x_1\backslash C') \arrow[r, "\pi_0(V)"] \arrow[d, dotted] & \pi_0(x_0\backslash C') \arrow[d, dotted] & \text{commutative.} \\
F(x_1) \arrow[r, "F(v)"] & F(x_0) &
\end{tikzcd}
\]
We consider the presheaf of sets
\[ x \longmapsto \pi_0(x\backslash C') \qquad \text{on } C, \]
let us call it $\pi_0(C'/C)$, we will have a canonical morphism
\[ \boxed{\varprojlim_{C'} \alpha^*(F) \xrightarrow{\ \sim\ } \operatorname{Hom}_{\substack{\text{presheaves}\\ \text{of sets on $C$}}} (\pi_0(C'/C), F),} \]
probably always \texttt{[an]} isomorphism.

If we work in the derived category of presheaves of $k$-modules ($k$ a ring), the preceding formula suggests conjecturally
\[ \boxed{\mathbf{R}\Gamma(C', \alpha^*F) \simeq \mathbf{R}\underbrace{\operatorname{Hom}}_{\text{internal Hom}}\ (\underbrace{\mathbf{L}\alpha_!(k_{C'})}_{\substack{\text{relative}\\ \text{homology of}\\ \text{$C'$ over $C$}}}, F).} \]

\msp{page 8}
The canonical augmentation
\[ \mathbf{L}\alpha_!(k_{C'}) \longrightarrow k_C \]
gives back, by taking the $\mathbf{R}\mathrm{Hom}$ into $F$,
\[ \mathbf{R}\Gamma(C,F) \longrightarrow \mathbf{R}\operatorname{Hom}(\mathbf{L}\alpha_!(k_{C'}), F) \xleftarrow{\ \sim\ } \mathbf{R}\Gamma(C', \alpha^*F) \]
the canonical homomorphism of ``contravariance'' of cohomology.

If on the other hand we work with the nonadditive derivator HOT, thus $F$ a ``relative homotopy type'' over $C$, we must have
\[ \boxed{\mathbf{R}\Gamma_{\mathrm{HOT}}(C', \alpha^*F) \xrightarrow{\ \sim\ } \mathbf{R}\Gamma(C\,,\ \underbrace{\operatorname{Hom}}_{\substack{\text{internal Hom}\\ \text{of the category of}\\ \text{relative homotopy}\\ \text{types over $C$}}}\ (\ \alpha_!\ (\ \underbrace{e_{C'}}_{\substack{\text{final object}\\ \text{of HOT($C'$),}\\ \text{final relative}\\ \text{homotopy type}\\ \text{over $C'$}}}\ ),\ F)).} \]

\printedpage{6}
% head: A. Grothendieck : Dérivateurs, Ch. II, Édition des Universités de Montpellier II et Paris VII
% foot: Version 8 juillet 2001
In fact, these formulas are quasi-tautological, if we admit the existence of the functors $\alpha_!$ as \textit{left} adjoints to the $\alpha^*$. In the abelian case, say,
\msp{page 9}
the isomorphy of the $H^0$ of the two sides is the adjunction formula
\[ \operatorname{Hom}_{C'}(k_{C'}, \alpha^*F) \simeq \operatorname{Hom}_C(\alpha_!(k_{C'}), F), \]
whence the isomorphy for the $H^i$ of the two sides, by shifting. Likewise in the case of HOT, the canonical isomorphy of the $\pi_0$ of the two sides is the adjunction formula. To have the finer isomorphism in $\mathbf{D}(e)$, the fundamental category of the derivator considered, in terms of internal Hom, we must admit the existence of unit objects (the $k_X$ in one case, the $e_X$ \texttt{[rather} $\varepsilon_X$\texttt{]} in the other), compatible with inverse images, in such a way \texttt{[that]} $\Homint(\varepsilon_X, \underbrace{F}_{\in\mathbf{D}_X}) \simeq F$. We will have
\[ \Gamma(C', \alpha^*F) \overset{\text{def}}{=} f'_*(\alpha^*(F)) \simeq f_*(\alpha_*\alpha^*(F)), \]
and the question is therefore to give an interpretation of $\alpha_*\alpha^*F$ as
\[ \alpha_*\alpha^*F \simeq \Homint(\alpha_!(\varepsilon_{C'}), F) \qquad \text{in } \mathbf{D}(C). \]
For this, we are going to write
\[ \alpha^*(F) \simeq \Homint(\varepsilon_{C'}, F), \]
and the formula must follow from an ``internal adjunction formula''
\msp{page 10}
\[ ? \qquad\qquad \boxed{\alpha_*(\Homint_{C'}(G', \alpha^*F)) \simeq \Homint_C(\alpha_!G', F).} \]
(\textit{``Internal'' adjunction formula,} for $\alpha_!$ and $\alpha^*$, valid in principle for \textit{every} $\alpha : C' \longrightarrow C$ in Diag, every $G'$ in $\mathbf{D}(C')$ and every $F$ in $\mathbf{D}(C)$.)

How to prove this formula? By taking the product operation $\otimes$ associated with the operation Hom\,int, we are reduced (by taking the associated contravariant functors) to a formula
\[ G \otimes \alpha_!G' \xrightarrow{\ \sim\ } \alpha_!(\alpha^*G \otimes G'). \]
We find by adjunction a canonical homomorphism, we must see whether it is always \texttt{[an]} isomorphism.

\printedpage{7}
% head: A. Grothendieck : Dérivateurs, Ch. II, Édition des Universités de Montpellier II et Paris VII
% foot: Version 8 juillet 2001
\msp{page 11}

\section*{2\quad Locally filtering categories are the locally totally $0$-connected categories}

Let $\mathcal{T}$ be a topos. We say that $\mathcal{T}$ is \textit{irreducible} if the ordered set $\operatorname{Ouv}(\mathcal{T})$ of subobjects of $e_{\mathcal{T}}$ has the following property: a) $\operatorname{Ouv}(\mathcal{T})$ is not reduced to one element, i.e. $\emptyset_{\mathcal{T}} \longrightarrow e_{\mathcal{T}}$ is not an isomorphism ($\mathcal{T}$ is not identical to its empty subtopos). b) If $U'$, $U''$ are two ``nonempty'' objects, i.e. noninitial, of $\operatorname{Ouv}(\mathcal{T})$, their intersection is ``nonempty''. It amounts to the same for b) to say that if $\mathcal{T}'$, $\mathcal{T}''$ are two closed subtopoi distinct from $\mathcal{T}$, then their supremum $\mathcal{T}' \cup \mathcal{T}''$ is distinct from $\mathcal{T}$, or again that every open of $\mathcal{T}$ is connected. We say that $\mathcal{T}$ is \textit{locally irreducible} if a) $\mathcal{T}$ is locally connected (i.e. each sheaf on $\mathcal{T}$ is a sum of 0-connected sheaves), and if b) for every 0-connected sheaf $U$ on $\mathcal{T}$, the induced topos $\mathcal{T}/U$ is irreducible (i.e. the ordered set $\operatorname{Ouv}(U)$ of subobjects of $U$ has property b) above ( a) simply means that $U \not\simeq \emptyset_{\mathcal{T}}$, i.e. that $U$ is not an initial object in $\mathcal{T}^\wedge$; or more simply): every subsheaf of a connected sheaf is connected (\footnote[5]{\textbf{NB} Irreducible does not imply locally irreducible, any more than connected does not imply locally connected etc.}) (\footnote[6]{$\mathcal{T}$ is locally irreducible if and only if it has a set of generators $(U_\alpha)$, with the $U_\alpha$ irreducible (hence 0-connected, hence $\mathcal{T}$ is locally connected).}).

When $\mathcal{T}$ is of the form $\operatorname{Top}(T)$, where $T$ is a topological space, then $\mathcal{T}$ irreducible means that $T$ is irreducible (for then $\operatorname{Ouv}(\mathcal{T}) \simeq \operatorname{Ouv}(T)$). This already implies trivially that $T$ is locally connected (every open is a disjoint sum of connected opens), thus $\operatorname{Top}(\mathcal{T})$ is locally connected. Moreover, this also implies that $\operatorname{Top}(T)$ is locally irreducible (\footnote[7]{make the proof explicit.}). (\textbf{NB} The 0-connected objects % [sic] \operatorname{Top}(\mathcal{T}) -> \operatorname{Top}(T), Top s'applique à l'espace T et non au topos \mathcal{T}
\msp{page 12}
of $\operatorname{Top}(T)^\wedge$, i.e. the 0-connected sheaves of sets on $T$, are those whose ``generic fiber'' has exactly one point \dots) Thus $\mathcal{T} = \operatorname{Top}(T)$ is locally irreducible if and only if $T$ \texttt{[is]} locally connected, and the connected components of $T$ are irreducible, i.e. if and only if $T = \coprod_i T_i$, the $T_i$ irreducible.

If $\mathcal{T}$ is of the form $\operatorname{Top}(X)$, $X$ a category, we know that $\mathcal{T}$ is locally connected (and even locally $\infty$-connected). We see on the other hand that $\mathcal{T}$ is irreducible if and only if $\operatorname{Esp}(X)$ (the associated topological space, whose underlying set is $\operatorname{Ob} X$) is irreducible (\footnote[8]{for the opens of $\operatorname{Top}(X)$, those of $\operatorname{Esp}(X)$ and those of $X$ correspond to one another.}); and we then say that $X$ is irreducible (\footnote[9]{This also means that for $x$, $y$ in $X$, $\exists\, z$ with \begin{tikzcd}[ampersand replacement=\&, row sep=0.2em, column sep=1.2em, baseline=(\tikzcdmatrixname-2-1.base)] \& x \\ z \arrow[ur] \arrow[dr] \& \\ \& y \end{tikzcd}, i.e. $\operatorname{Ord}(X)$ is left filtering.}). As the connected components of $X$ and those of $\operatorname{Esp}(X)$ or $\operatorname{Top}(X) = \mathcal{T}$ correspond to one another, we see therefore that the following conditions are equivalent: a) the topos $\operatorname{Top}(X)$ is locally irreducible, b)
\printedpage{8}
% head: A. Grothendieck : Dérivateurs, Ch. II, Édition des Universités de Montpellier II et Paris VII
% foot: Version 8 juillet 2001
$\operatorname{Esp}(X)$ is locally irreducible, c) the connected components of $X$ are irreducible, d) the $X/x$ \texttt{[are]} irreducible. We then say that $X$ is locally irreducible (\footnote[10]{False, we have (a $\Longleftrightarrow$ d) $\Longrightarrow$ (b $\Longleftrightarrow$ c) $X$ irreducible does \textit{not} imply that the $X/x$ are so, example $X = (0 \overset{\alpha}{\underset{\beta}{\rightrightarrows}} 1)$ $(\alpha \neq \beta)$, we have $X/1 =$ \begin{tikzcd}[ampersand replacement=\&, row sep=0.2em, column sep=1.2em, baseline=(\tikzcdmatrixname-2-2.base)] 0' \arrow[dr] \& \\ \& 1 \\ 0'' \arrow[ur] \& \end{tikzcd}, which is \textit{not} irreducible, nor locally irreducible since it is connected. Yet it is étale over $X$.}).

This therefore implies that for every $X'$ étale over $X$, with $X'$ 0-connected, $X'$ is irreducible.

\textit{State equivalent conditions.}
\begin{enumerate}
\item[a)] $\operatorname{Top}(X)$ is locally irreducible.
\item[b)] The $X/x$ \texttt{[are]} irreducible.
\item[c)] Every 0-connected object of $X^\wedge$ is irreducible.
\item[d)] The connected components of $X$ are weakly cofiltering (i.e. weakly left filtering).
\item[e)] Condition PS $1_{\mathrm{g}}$ of SGA 4, I 2.
\end{enumerate}
If $X$ is \texttt{[a]} category, $X$ irreducible also means that $X \neq \emptyset$ and $\forall\, x, y \in \operatorname{Ob} X$, \texttt{[there]} exists \texttt{[an]} $s$ with \begin{tikzcd}[row sep=0.2em, column sep=1.2em, baseline=(\tikzcdmatrixname-2-1.base)] & x \\ s \arrow[ur] \arrow[dr] & \\ & y \end{tikzcd}. (This also means that $X$ is \textit{weakly filtering} (\footnote[11]{It is better to call it \textit{pseudo-filtering,} and let us reserve weakly filtering for: locally weakly filtering + connected.}) on the left, or weakly cofiltering.) $X$ locally irreducible means that for $x$, $y$ \textit{belonging to the same connected component,} \texttt{[there]} exists \texttt{[an]} $s$ with \begin{tikzcd}[row sep=0.2em, column sep=1.2em, baseline=(\tikzcdmatrixname-2-1.base)] & x \\ s \arrow[ur] \arrow[dr] & \\ & y \end{tikzcd}, % [sic] irreductible -> irréductible, accent omis
\msp{page 13}
i.e. $X/[x] \cap X/[y] \neq \emptyset$ (where $[x]$ is regarded as an element of the ordered set $\operatorname{Ob} X$). It again amounts to the same to say that the induced categories $X/x$ are irreducible (\footnote[12]{false, cf. marginal annotation above.}).

\textbf{Example.} A category \textit{having an initial object,} or merely a ``smallest object'', is irreducible. Conversely, if the \textit{ordered} set $\operatorname{Ord}(X)$ associated with $X$ (quotient of $\operatorname{Ob}(X)$ by the equivalence relation associated with its preorder) is finite, e.g. $X$ finite, or only locally finite, e.g. $X$ locally finite (i.e. the categories $X/x$ finite), then if $X$ is irreducible, $\operatorname{Ob} X$ has a smallest element, i.e. $X$ has a smallest object. This is true more generally still, if $\operatorname{Ord}(X)$ is artinian, i.e. every decreasing sequence in $\operatorname{Ord}(X)$ is stationary. % [sic] initiale -> initial, accord avec « objet »

But the interesting condition is the topos-theoretic condition of \textit{local} irreducibility, which is equivalent to: $I$ left pseudo-filtering, i.e. the connected components of $I$ filtering, i.e. conditions PS 1, PS 2

PS 1\enspace Every diagram \begin{tikzcd}[row sep=0.2em, column sep=1.2em, baseline=(\tikzcdmatrixname-2-2.base)] x \arrow[dr] & \\ & s \\ y \arrow[ur] & \end{tikzcd} can be completed to \begin{tikzcd}[row sep=0.2em, column sep=1.2em, baseline=(\tikzcdmatrixname-2-1.base)] & x \arrow[dr] & \\ z \arrow[ur] \arrow[dr] & & s \\ & y \arrow[ur] & \end{tikzcd} \textit{commutative.}

PS 2\enspace For every double arrow $x \rightrightarrows y$, \texttt{[there]} exists an equalizer $z \longrightarrow x \rightrightarrows y$.

\end{document}
